% GUID v6, Primer P6 (label box:cosmo): cosmology. Background only; it states no result of the paper.
% Place it at the start of Section sec:cosmo (marker %PRIMER:cosmo in part1.tex): % GUID v6, Primer P6 (label box:cosmo): cosmology. Background only; it states no result of the paper.
% Place it at the start of Section sec:cosmo (marker %PRIMER:cosmo in part1.tex): % GUID v6, Primer P6 (label box:cosmo): cosmology. Background only; it states no result of the paper.
% Place it at the start of Section sec:cosmo (marker %PRIMER:cosmo in part1.tex): % GUID v6, Primer P6 (label box:cosmo): cosmology. Background only; it states no result of the paper.
% Place it at the start of Section sec:cosmo (marker %PRIMER:cosmo in part1.tex): \input{primers/p6-cosmo}
\begin{primer}{Cosmology}{box:cosmo}
\emph{The expanding universe.} On large scales the universe is homogeneous and isotropic. Distances grow with a scale factor $a(t)$, light
from the past is redshifted by $1+z=a_0/a$, and the expansion rate $H=\dot a/a$ obeys, for a flat universe,
$H^2=8\pi G\rho/3+\Lambda c^2/3$. Each component's share is $\Omega_i=\rho_i/\rho_c$ with the critical density $\rho_c=3H^2/8\pi G$; the
value today is $H_0$.

\emph{$\Lambda$CDM.} The standard model of cosmology adds to the known particles a cold, collisionless \emph{dark matter} and a
cosmological constant $\Lambda$. Six parameters fix it: the densities of baryons and dark matter ($\Omega_bh^2$, $\Omega_ch^2$), the angular
size of the sound horizon, the optical depth to reionisation, and the amplitude $A_s$ and tilt $n_s$ of the primordial
fluctuations~\cite{planck18}.

\emph{Inflation.} An early epoch of nearly exponential expansion, $a\propto e^{Ht}$, flattens and smooths the universe and stretches
quantum fluctuations into the seeds of galaxies~\cite{guth}. It is driven by a scalar field, the \emph{inflaton}, with potential $V$, and
its length is counted in \emph{e-folds}, $N=\ln(a_{\rm end}/a_{\rm start})$. In \emph{slow roll} the field creeps down a flat potential;
with the reduced Planck mass $M_{\rm pl}$, the parameters $\epsilon=\tfrac12M_{\rm pl}^2(V'/V)^2$ and $\eta_V=M_{\rm pl}^2V''/V$ are small. The
fluctuations then have a nearly scale-invariant spectrum with tilt $n_s\approx1-6\epsilon+2\eta_V$, and gravitational waves with a
tensor-to-scalar ratio $r\approx16\epsilon$~\cite{planck2020x}.

\emph{Reheating and oscillons.} After inflation the inflaton oscillates about the minimum of its potential and decays into ordinary
particles, which \emph{reheats} the universe~\cite{kls}. In a broad class of models it first fragments into \emph{oscillons}, long-lived
localised lumps of the field~\cite{amin}.

\emph{The probes.} The cosmic microwave background (CMB) was released at recombination, $z\approx1100$. Its acoustic peaks are the imprint of
sound waves in the early plasma: their spacing fixes the sound horizon and their heights the densities~\cite{planck18}. The same sound
horizon, about 150\,Mpc today, is imprinted on the clustering of galaxies as \emph{baryon acoustic oscillations} (BAO), a standard ruler
for distance against redshift~\cite{desi2}.

\emph{Dark energy.} A cosmological constant has the equation of state $w=p/\rho c^2=-1$ at all times. DESI fits the
Chevallier--Polarski--Linder (CPL) form $w(a)=w_0+w_a(1-a)$~\cite{chevallier2001,linder2003}. A \emph{thawing} dark energy has $w_0>-1$
and $w_a<0$: it was close to $-1$ early and weakens today.

\emph{Data names and tensions.} P-ACT combines the CMB spectra of Planck and of the Atacama Cosmology Telescope (ACT, data release 6);
P-ACT-LB adds CMB lensing (L) and DESI BAO (B)~\cite{louis2025}. The \emph{$H_0$ tension} is the disagreement between the Hubble constant
inferred from the CMB with $\Lambda$CDM, near 67\,km\,s$^{-1}$\,Mpc$^{-1}$, and the value from a distance ladder of Cepheids and supernovae,
near 73~\cite{planck18,riess2022}. A ladder calibrated on the tip of the red-giant branch lies in between~\cite{freedman2025}.
\end{primer}

\begin{primer}{Cosmology}{box:cosmo}
\emph{The expanding universe.} On large scales the universe is homogeneous and isotropic. Distances grow with a scale factor $a(t)$, light
from the past is redshifted by $1+z=a_0/a$, and the expansion rate $H=\dot a/a$ obeys, for a flat universe,
$H^2=8\pi G\rho/3+\Lambda c^2/3$. Each component's share is $\Omega_i=\rho_i/\rho_c$ with the critical density $\rho_c=3H^2/8\pi G$; the
value today is $H_0$.

\emph{$\Lambda$CDM.} The standard model of cosmology adds to the known particles a cold, collisionless \emph{dark matter} and a
cosmological constant $\Lambda$. Six parameters fix it: the densities of baryons and dark matter ($\Omega_bh^2$, $\Omega_ch^2$), the angular
size of the sound horizon, the optical depth to reionisation, and the amplitude $A_s$ and tilt $n_s$ of the primordial
fluctuations~\cite{planck18}.

\emph{Inflation.} An early epoch of nearly exponential expansion, $a\propto e^{Ht}$, flattens and smooths the universe and stretches
quantum fluctuations into the seeds of galaxies~\cite{guth}. It is driven by a scalar field, the \emph{inflaton}, with potential $V$, and
its length is counted in \emph{e-folds}, $N=\ln(a_{\rm end}/a_{\rm start})$. In \emph{slow roll} the field creeps down a flat potential;
with the reduced Planck mass $M_{\rm pl}$, the parameters $\epsilon=\tfrac12M_{\rm pl}^2(V'/V)^2$ and $\eta_V=M_{\rm pl}^2V''/V$ are small. The
fluctuations then have a nearly scale-invariant spectrum with tilt $n_s\approx1-6\epsilon+2\eta_V$, and gravitational waves with a
tensor-to-scalar ratio $r\approx16\epsilon$~\cite{planck2020x}.

\emph{Reheating and oscillons.} After inflation the inflaton oscillates about the minimum of its potential and decays into ordinary
particles, which \emph{reheats} the universe~\cite{kls}. In a broad class of models it first fragments into \emph{oscillons}, long-lived
localised lumps of the field~\cite{amin}.

\emph{The probes.} The cosmic microwave background (CMB) was released at recombination, $z\approx1100$. Its acoustic peaks are the imprint of
sound waves in the early plasma: their spacing fixes the sound horizon and their heights the densities~\cite{planck18}. The same sound
horizon, about 150\,Mpc today, is imprinted on the clustering of galaxies as \emph{baryon acoustic oscillations} (BAO), a standard ruler
for distance against redshift~\cite{desi2}.

\emph{Dark energy.} A cosmological constant has the equation of state $w=p/\rho c^2=-1$ at all times. DESI fits the
Chevallier--Polarski--Linder (CPL) form $w(a)=w_0+w_a(1-a)$~\cite{chevallier2001,linder2003}. A \emph{thawing} dark energy has $w_0>-1$
and $w_a<0$: it was close to $-1$ early and weakens today.

\emph{Data names and tensions.} P-ACT combines the CMB spectra of Planck and of the Atacama Cosmology Telescope (ACT, data release 6);
P-ACT-LB adds CMB lensing (L) and DESI BAO (B)~\cite{louis2025}. The \emph{$H_0$ tension} is the disagreement between the Hubble constant
inferred from the CMB with $\Lambda$CDM, near 67\,km\,s$^{-1}$\,Mpc$^{-1}$, and the value from a distance ladder of Cepheids and supernovae,
near 73~\cite{planck18,riess2022}. A ladder calibrated on the tip of the red-giant branch lies in between~\cite{freedman2025}.
\end{primer}

\begin{primer}{Cosmology}{box:cosmo}
\emph{The expanding universe.} On large scales the universe is homogeneous and isotropic. Distances grow with a scale factor $a(t)$, light
from the past is redshifted by $1+z=a_0/a$, and the expansion rate $H=\dot a/a$ obeys, for a flat universe,
$H^2=8\pi G\rho/3+\Lambda c^2/3$. Each component's share is $\Omega_i=\rho_i/\rho_c$ with the critical density $\rho_c=3H^2/8\pi G$; the
value today is $H_0$.

\emph{$\Lambda$CDM.} The standard model of cosmology adds to the known particles a cold, collisionless \emph{dark matter} and a
cosmological constant $\Lambda$. Six parameters fix it: the densities of baryons and dark matter ($\Omega_bh^2$, $\Omega_ch^2$), the angular
size of the sound horizon, the optical depth to reionisation, and the amplitude $A_s$ and tilt $n_s$ of the primordial
fluctuations~\cite{planck18}.

\emph{Inflation.} An early epoch of nearly exponential expansion, $a\propto e^{Ht}$, flattens and smooths the universe and stretches
quantum fluctuations into the seeds of galaxies~\cite{guth}. It is driven by a scalar field, the \emph{inflaton}, with potential $V$, and
its length is counted in \emph{e-folds}, $N=\ln(a_{\rm end}/a_{\rm start})$. In \emph{slow roll} the field creeps down a flat potential;
with the reduced Planck mass $M_{\rm pl}$, the parameters $\epsilon=\tfrac12M_{\rm pl}^2(V'/V)^2$ and $\eta_V=M_{\rm pl}^2V''/V$ are small. The
fluctuations then have a nearly scale-invariant spectrum with tilt $n_s\approx1-6\epsilon+2\eta_V$, and gravitational waves with a
tensor-to-scalar ratio $r\approx16\epsilon$~\cite{planck2020x}.

\emph{Reheating and oscillons.} After inflation the inflaton oscillates about the minimum of its potential and decays into ordinary
particles, which \emph{reheats} the universe~\cite{kls}. In a broad class of models it first fragments into \emph{oscillons}, long-lived
localised lumps of the field~\cite{amin}.

\emph{The probes.} The cosmic microwave background (CMB) was released at recombination, $z\approx1100$. Its acoustic peaks are the imprint of
sound waves in the early plasma: their spacing fixes the sound horizon and their heights the densities~\cite{planck18}. The same sound
horizon, about 150\,Mpc today, is imprinted on the clustering of galaxies as \emph{baryon acoustic oscillations} (BAO), a standard ruler
for distance against redshift~\cite{desi2}.

\emph{Dark energy.} A cosmological constant has the equation of state $w=p/\rho c^2=-1$ at all times. DESI fits the
Chevallier--Polarski--Linder (CPL) form $w(a)=w_0+w_a(1-a)$~\cite{chevallier2001,linder2003}. A \emph{thawing} dark energy has $w_0>-1$
and $w_a<0$: it was close to $-1$ early and weakens today.

\emph{Data names and tensions.} P-ACT combines the CMB spectra of Planck and of the Atacama Cosmology Telescope (ACT, data release 6);
P-ACT-LB adds CMB lensing (L) and DESI BAO (B)~\cite{louis2025}. The \emph{$H_0$ tension} is the disagreement between the Hubble constant
inferred from the CMB with $\Lambda$CDM, near 67\,km\,s$^{-1}$\,Mpc$^{-1}$, and the value from a distance ladder of Cepheids and supernovae,
near 73~\cite{planck18,riess2022}. A ladder calibrated on the tip of the red-giant branch lies in between~\cite{freedman2025}.
\end{primer}

\begin{primer}{Cosmology}{box:cosmo}
\emph{The expanding universe.} On large scales the universe is homogeneous and isotropic. Distances grow with a scale factor $a(t)$, light
from the past is redshifted by $1+z=a_0/a$, and the expansion rate $H=\dot a/a$ obeys, for a flat universe,
$H^2=8\pi G\rho/3+\Lambda c^2/3$. Each component's share is $\Omega_i=\rho_i/\rho_c$ with the critical density $\rho_c=3H^2/8\pi G$; the
value today is $H_0$.

\emph{$\Lambda$CDM.} The standard model of cosmology adds to the known particles a cold, collisionless \emph{dark matter} and a
cosmological constant $\Lambda$. Six parameters fix it: the densities of baryons and dark matter ($\Omega_bh^2$, $\Omega_ch^2$), the angular
size of the sound horizon, the optical depth to reionisation, and the amplitude $A_s$ and tilt $n_s$ of the primordial
fluctuations~\cite{planck18}.

\emph{Inflation.} An early epoch of nearly exponential expansion, $a\propto e^{Ht}$, flattens and smooths the universe and stretches
quantum fluctuations into the seeds of galaxies~\cite{guth}. It is driven by a scalar field, the \emph{inflaton}, with potential $V$, and
its length is counted in \emph{e-folds}, $N=\ln(a_{\rm end}/a_{\rm start})$. In \emph{slow roll} the field creeps down a flat potential;
with the reduced Planck mass $M_{\rm pl}$, the parameters $\epsilon=\tfrac12M_{\rm pl}^2(V'/V)^2$ and $\eta_V=M_{\rm pl}^2V''/V$ are small. The
fluctuations then have a nearly scale-invariant spectrum with tilt $n_s\approx1-6\epsilon+2\eta_V$, and gravitational waves with a
tensor-to-scalar ratio $r\approx16\epsilon$~\cite{planck2020x}.

\emph{Reheating and oscillons.} After inflation the inflaton oscillates about the minimum of its potential and decays into ordinary
particles, which \emph{reheats} the universe~\cite{kls}. In a broad class of models it first fragments into \emph{oscillons}, long-lived
localised lumps of the field~\cite{amin}.

\emph{The probes.} The cosmic microwave background (CMB) was released at recombination, $z\approx1100$. Its acoustic peaks are the imprint of
sound waves in the early plasma: their spacing fixes the sound horizon and their heights the densities~\cite{planck18}. The same sound
horizon, about 150\,Mpc today, is imprinted on the clustering of galaxies as \emph{baryon acoustic oscillations} (BAO), a standard ruler
for distance against redshift~\cite{desi2}.

\emph{Dark energy.} A cosmological constant has the equation of state $w=p/\rho c^2=-1$ at all times. DESI fits the
Chevallier--Polarski--Linder (CPL) form $w(a)=w_0+w_a(1-a)$~\cite{chevallier2001,linder2003}. A \emph{thawing} dark energy has $w_0>-1$
and $w_a<0$: it was close to $-1$ early and weakens today.

\emph{Data names and tensions.} P-ACT combines the CMB spectra of Planck and of the Atacama Cosmology Telescope (ACT, data release 6);
P-ACT-LB adds CMB lensing (L) and DESI BAO (B)~\cite{louis2025}. The \emph{$H_0$ tension} is the disagreement between the Hubble constant
inferred from the CMB with $\Lambda$CDM, near 67\,km\,s$^{-1}$\,Mpc$^{-1}$, and the value from a distance ladder of Cepheids and supernovae,
near 73~\cite{planck18,riess2022}. A ladder calibrated on the tip of the red-giant branch lies in between~\cite{freedman2025}.
\end{primer}
